\chapter{Fehlerkorpus} % expect: STR-ORIENTIERUNG
\label{chap:fehler}

\section{Typografie}
\label{sec:typo}

Dieser Abschnitt enthält typografische Fehler in kurzen Sätzen.
Die Stränge laufen parallel; die Fusion folgt danach. % expect: TYP-SEMIKOLON
Ein System — das adaptiert — wurde gebaut. % expect: TYP-GEDANKENSTRICH
Die Adaption -- so die Annahme -- wirkt. % expect: TYP-HALBGEVIERT
Das Phänomen heißt "Persona Collapse" in der Fachsprache. % expect: TYP-ANFUEHRUNG
Das betrifft z.B. die Stimme. % expect: TYP-ABKUERZUNG
Der Anteil liegt bei 15\% der Beiträge. % expect: TYP-PROZENT
Die Studie umfasst N = 40 Personen. % expect: TYP-EINHEIT
Die Studie umfasst 3 Phasen. % expect: TYP-ZAHLWORT

\begin{itemize}
    \item Erster Punkt der Liste.

    \item Zweiter Punkt der Liste. % expect: TYP-AUFZAEHLUNG-LEER
\end{itemize}

\section{Sprache}
\label{sec:sprache}

Dieser Abschnitt enthält sprachliche Fehler in kurzen Sätzen.
Wir zeigen, dass die Adaption wirkt. % expect: SPR-ICH-WIR
Das System liefert beeindruckende Ergebnisse. % expect: SPR-WERTUNG
Natürlich beeinflusst die Stimme die Wahrnehmung. % expect: SPR-SELBSTVERSTAENDLICH
Wie bereits erwähnt, werden die Stränge gewichtet. % expect: SPR-FLOSKEL
Dies könnte möglicherweise zutreffen. % expect: SPR-HEDGING
Die Fusion ist eigentlich ein Mittelwert. % expect: SPR-FUELLWORT
Zuerst zunächst wird der Selbsttest bearbeitet. % expect: SPR-DOPPELUNG
Die Stimme hat einen signifikanten Einfluss auf das Erleben. % expect: SPR-SIGNIFIKANT
Die Durchführung der Erhebung erfolgt nach Abschluss der Implementierung, Prüfung, Dokumentation und Konfiguration. % expect: SPR-NOMINALSTIL

\section{Satzbau}
\label{sec:satz}

Dieser Abschnitt enthält Fehler im Satzbau.
Der Prototyp verbindet die Erkennung der Persönlichkeit aus der Sprache mit einer dynamischen Anpassung des Agenten an die erkannte Persönlichkeit und gibt diese Anpassung zusätzlich über die Stimme aus, wobei alle Schritte in einer gemeinsamen Verarbeitungskette ablaufen, die für jede Äußerung neu startet. % expect: SAT-LAENGE
Die Arbeit zeigt, dass der Agent, der die Werte, die der Strang liefert, nutzt, besser abschneidet. % expect: SAT-SCHACHTEL
Die Ergebnisse der Evaluation der Nutzerstudie des Systems werden berichtet. % expect: SAT-GENITIV

\section{Begriffe}
\label{sec:begriffe}

Dieser Abschnitt enthält Fehler bei Begriffen.
Jeder Werkzeugaufruf wird gespeichert. % expect: TERM-VERBOTEN
Der Nutzer spricht mit dem Agenten.
Der Nutzer bewertet den Agenten.
Der Anwender füllt den Fragebogen aus. % expect: TERM-VARIANTEN
Ein KI Voice Agent passt sich an. % expect: TERM-KOMPOSITUM, TERM-ABKUERZUNG
Das LLM erzeugt die Antwort. % expect: TERM-ABKUERZUNG
Das \gls{a:LLM} erzeugt die Antwort.
Der Begriff \emph{Persona Collapse} beschreibt den Identitätsverlust.
Ein \emph{Persona Collapse} tritt bei starker Anpassung auf. % expect: TERM-EMPH

\section{Zitation}
\label{sec:zitation}

Dieser Abschnitt enthält Fehler bei Zitaten.
Die Studie von \textcite{unbekannt2020} zeigt einen Effekt. % expect: ZIT-SCHLUESSEL
Peters et al. (2024) zeigen einen Effekt. % expect: ZIT-MANUELL
Studien zeigen, dass Extraversion hörbar ist. % expect: ZIT-BELEG
Die Persönlichkeit gilt als stabil \cite{goldberg1990}. % expect: ZIT-BEFEHL
Die Dokumentation steht unter \url{https://example.org}. % expect: ZIT-URL
Goldberg beschreibt es als \enquote{eine robuste Struktur aus fünf breiten Faktoren der Persönlichkeit} \parencite{goldberg1990}. % expect: ZIT-WOERTLICH

\section{Formatierung}
\label{sec:format}

Dieser Abschnitt enthält Fehler bei der Formatierung.
Die Adaption ist \textbf{entscheidend} für das Ergebnis. % expect: FMT-FETT
Abbildung \ref{fig:ohnequelle} zeigt das System. % expect: FMT-REF
Der Verweis auf \cref{fig:gibtesnicht} läuft ins Leere. % expect: FMT-VERWEIS-UNBEKANNT

\begin{figure}[htbp] % expect: FMT-ABB-QUELLE
    \centering
    \rule{2cm}{1cm}
    \caption{Architektur ohne Herkunft}
    \label{fig:ohnequelle}
\end{figure}

\begin{table}[htbp] % expect: FMT-CAPTION, FMT-LABEL
    \centering
    \begin{tabular}{|l|r|} % expect: FMT-VLINIE
        \hline % expect: FMT-VLINIE
        A & 1 \\
    \end{tabular}
\end{table}

\begin{figure}[htbp] % expect: FMT-VERWEIS-FEHLT
    \centering
    \rule{2cm}{1cm}
    \caption{Nie referenziert. Eigene Darstellung}
    \label{fig:nie}
\end{figure}

\begin{figure}[htbp] % expect: FMT-VERWEIS-REIHENFOLGE
    \centering
    \rule{2cm}{1cm}
    \caption{Erst danach referenziert. Eigene Darstellung}
    \label{fig:spaet}
\end{figure}

Erst nach der Abbildung folgt der Verweis auf \cref{fig:spaet}.
Hier stehen zwei Sätze in einer Zeile. Das ist der zweite Satz. % expect: FMT-QUELLTEXT-ZEILE

\section{Marker}
\label{sec:marker}

Dieser Abschnitt enthält Marker der schreibenden Person.
Die Übersetzung verändert die Kennwerte kaum \quellefehlt{Windsor et al. 2019}. % expect: MARK-QUELLE-FEHLT
\offen{Zahl der Ausschlüsse ergänzen.} % expect: MARK-OFFEN
Die Stimme bleibt konstant.
Das Ergebnis ist beeindruckend. % thesis-editor: ignore SPR-WERTUNG

\section{Struktur} % expect: STR-ORIENTIERUNG
\label{sec:struktur}

\subsection{Ohne Einleitung}

Dieser Unterabschnitt hat genug Text.
Er besteht aus drei Sätzen.
Damit ist er lang genug.

\subsection{Zu kurz} % expect: STR-KURZ

Nur ein Satz steht hier.

\subsection{Formel}
\label{subsec:formel}

Die Fusion bildet einen gewichteten Mittelwert.
Die folgende Formel beschreibt ihn.
\begin{equation} % expect: FORM-ERKLAERUNG
    s = w_1 x_1 + w_2 x_2
\end{equation}
Die Gewichte sind fest.
